\documentclass[12pt]{article}
\usepackage[a4paper,margin=1in]{geometry}
\usepackage{amsmath,amssymb,mathtools}
\usepackage{setspace,enumitem,booktabs,lmodern,microtype}
\usepackage{titlesec}
\setstretch{1.25}
\setlength{\parindent}{0pt}
\setlength{\parskip}{6pt}

\titleformat{\section}{\large\bfseries}{}{0em}{}
\titleformat{\subsection}{\normalsize\bfseries}{}{0em}{}

\begin{document}

\begin{center}
{\Large \textbf{Midterm Exam 2}} \\[0.9em]
\textbf{ECON 320 -- Introduction to Mathematical Economics} \\[0.7em]

November 05, 2025
 \\[0.8em]
Instructor: Muhebullah Karimzada
\end{center}

\hrule

\textbf{Instructions:}
\begin{itemize}
\item You are allowed \textbf{75 minutes} to complete the exam.
    \item Show all work. Unsupported answers receive little credit.

    \item Unless stated otherwise, all variables are real and differentiable as needed.
\end{itemize}

\hrule

\section*{Problem 1 (12 points)}

A firm’s profit function is
\[
\pi(Q) = 120Q - 6Q^2 - 180.
\]

\begin{enumerate}[label=(\alph*),leftmargin=1.5em]
    \item Derive the first-order condition and find the critical point \(Q^*\). \hfill (4 pts)
    \item Compute the second derivative and determine whether \(Q^*\) is a maximum or minimum. \hfill (3 pts)
    \item Find the maximum profit \(\pi^*\). \hfill (3 pts)
    \item Briefly interpret \(Q^*\) and the curvature of \(\pi(Q)\) economically. \hfill (2 pts)
\end{enumerate}



\section*{Problem 2  (12 points)}

A population grows according to
\[
P(t) = 200 e^{0.03t},
\]
and a technology index grows as
\[
A(t) = e^{0.02t}.
\]
Output is given by a Cobb–Douglas form \(Y(t) = A(t) P(t)^{0.5}\).
\newpage
\begin{enumerate}[label=(\alph*),leftmargin=1.5em]
    \item Find the instantaneous growth rate of \(Y(t)\) by differentiating \(\ln Y(t)\). \hfill (4 pts)
    \item Compute \(Y(10)\) (rounded to two decimals). \hfill (3 pts)
    \item Suppose the population growth rate rises from 3\% to 4\%, while technology remains the same. Compute the new growth rate of \(Y(t)\). \hfill (3 pts)
    \item Interpret the difference in growth rates: what does it imply about the contribution of population versus technology? \hfill (2 pts)
\end{enumerate}



\section*{Problem 3 (13 points)}

Let
\[
f(x,y) = 12x + 8y - x^2 - y^2 - xy.
\]

\begin{enumerate}[label=(\alph*),leftmargin=1.5em]
    \item Compute \(f_x\) and \(f_y\) and find all stationary points. \hfill (4 pts)
    \item Form the Hessian matrix \(H\) and its determinant. \hfill (4 pts)
    \item Use the Hessian to classify the stationary point (maximum, minimum, or saddle). \hfill (3 pts)
    \item Provide an economic interpretation of the curvature of \(f(x,y)\). \hfill (2 pts)
\end{enumerate}



\section*{Problem 4 (13 points)}

A consumer maximizes
\[
U(x,y) = x^{0.5} y^{0.5}
\]
subject to the budget constraint
\[
2x + y = 100.
\]

\begin{enumerate}[label=(\alph*),leftmargin=1.5em]
    \item Form the Lagrangian and derive the first-order conditions. \hfill (4 pts)
    \item Solve for \(x^*, y^*, \lambda^*\). \hfill (5 pts)
    \item Verify (briefly) that the bordered Hessian condition confirms a maximum. \hfill (2 pts)
    \item Interpret the meaning of \(\lambda^*\) and use it to estimate the change in utility if income increases by 5. \hfill (2 pts)
\end{enumerate}

\vspace{10pt}
\hrule
\vspace{8pt}



\end{document}
